\clearpage
\begin{examplebox}
\textbf{\textcolor{CExample}{例 1：初中起步——先平方，再加 3}}
\begin{center}\begin{tikzpicture}[font=\small,>=stealth,line join=round]
\draw[gray!40,rounded corners=3pt] (0,0) rectangle (14.2,8.3);\node at (7.1,7.700000000000001) {例1：内偶先合并，外层加三不改变同高};
\begin{scope}[shift={(4.1,2.161111111111111)},x=1.25cm,y=0.5111111111111111cm]\draw[->,gray] (-2.4,0)--(2.4,0) node[right]{$x$};\draw[->,gray] (0,-1)--(0,8) node[above]{$y$};\draw[blue,thick,solid,domain=-2.1:2.1,samples=180] plot(\x,{(\x)^2+3});\fill[orange!85!black] (-2,7) circle(2.2pt);\fill[orange!85!black] (2,7) circle(2.2pt);\end{scope}\node[anchor=west,align=left,text width=5.0cm] at (8.4,4.0) {$F=x^2+3$\\[10pt]$g=x^2$，$f(u)=u+3$\\[10pt]$F(2)=F(-2)=7$\\[10pt]外层本身非奇非偶};
\node[text width=12.6cm,align=center] at (7.1,.65) {向上平移不破坏左右同高；这不是要求外层也具有奇偶性。};
\end{tikzpicture}\end{center}


\textbf{题目：} 判断 $F(x)=x^2+3$ 的奇偶性，并用 $x=2,-2$ 解释。\\
\textbf{拆成两层：} $g(x)=x^2$，$f(u)=u+3$。\\
\textbf{算一对数：} $F(2)=4+3=7$，$F(-2)=4+3=7$。\\
\textbf{再说明任意 $x$：} $(-x)^2=x^2$，所以 $(-x)^2+3=x^2+3$；整个函数为\textbf{偶函数}。\\
\textbf{画图记忆：} 它是 $y=x^2$ 整体向上平移 3 个单位，仍关于 $y$ 轴对称。

\medskip

\end{examplebox}
\clearpage
\begin{examplebox}
\textbf{\textcolor{CExample}{例 2：内奇外奇——乘 2 后立方}}
\begin{center}\begin{tikzpicture}[font=\small,>=stealth,line join=round]
\draw[gray!40,rounded corners=3pt] (0,0) rectangle (14.2,8.3);\node at (7.1,7.700000000000001) {例2：先乘二再立方，整体保持反号};
\begin{scope}[shift={(4.1,3.9499999999999997)},x=2.142857142857143cm,y=0.19166666666666665cm]\draw[->,gray] (-1.4,0)--(1.4,0) node[right]{$x$};\draw[->,gray] (0,-12)--(0,12) node[above]{$y$};\draw[blue,thick,solid,domain=-1.1:1.1,samples=180] plot(\x,{8*(\x)^3});\fill[orange!85!black] (-1,-8) circle(2.2pt);\fill[orange!85!black] (1,8) circle(2.2pt);\end{scope}\node[anchor=west,align=left,text width=5.0cm] at (8.4,4.0) {$F=(2x)^3$\\[10pt]内层 $2x$ 为奇\\[10pt]外层 $u^3$ 为奇\\[10pt]$F(1)=8,\ F(-1)=-8$};
\node[text width=12.6cm,align=center] at (7.1,.65) {标出的两个点帮助观察；任意输入的结论由正文等式验证。};
\end{tikzpicture}\end{center}


\textbf{题目：} 判断 $F(x)=(2x)^3$。\\
\textbf{拆成两层：} $g(x)=2x$（奇），$f(u)=u^3$（奇）。\\
\textbf{算一对数：} $F(1)=2^3=8$，$F(-1)=(-2)^3=-8$。\\
\textbf{写出证明：}
\[
F(-x)=(-2x)^3=-(2x)^3=-F(x).
\]
因此 $F$ 为\textbf{奇函数}。图像关于原点对称。

\medskip

\end{examplebox}
\clearpage
\begin{examplebox}
\textbf{\textcolor{CExample}{例 3：内奇外偶——乘 2 后平方}}
\begin{center}\begin{tikzpicture}[font=\small,>=stealth,line join=round]
\draw[gray!40,rounded corners=3pt] (0,0) rectangle (14.2,8.3);\node at (7.1,7.700000000000001) {例3：内层虽反号，外层平方把两路合并};
\begin{scope}[shift={(4.1,2.161111111111111)},x=1.875cm,y=0.5111111111111111cm]\draw[->,gray] (-1.6,0)--(1.6,0) node[right]{$x$};\draw[->,gray] (0,-1)--(0,8) node[above]{$y$};\draw[blue,thick,solid,domain=-1.3:1.3,samples=180] plot(\x,{4*(\x)^2});\fill[orange!85!black] (-1,4) circle(2.2pt);\fill[orange!85!black] (1,4) circle(2.2pt);\end{scope}\node[anchor=west,align=left,text width=5.0cm] at (8.4,4.0) {$F=(2x)^2$\\[10pt]$g(1)=2,\ g(-1)=-2$\\[10pt]平方以后都变成 $4$\\[10pt]内奇外偶，整体偶};
\node[text width=12.6cm,align=center] at (7.1,.65) {不能把“内层为奇”误记成“整体必奇”；还需要检查外层。};
\end{tikzpicture}\end{center}


\textbf{题目：} 判断 $F(x)=(2x)^2$。\\
\textbf{拆成两层：} 仍是 $g(x)=2x$（奇），但 $f(u)=u^2$（偶）。\\
\textbf{算一对数：} $F(1)=4$，$F(-1)=4$。\\
\textbf{写出证明：}
\[
F(-x)=(-2x)^2=(2x)^2=F(x).
\]
因此 $F$ 为\textbf{偶函数}。\textbf{对比例 2：里面完全一样，仅仅把外面“立方”换成“平方”，奇偶性就变了。}

\medskip

\end{examplebox}
\clearpage
\begin{examplebox}
\textbf{\textcolor{CExample}{例 4：看起来复杂，里面是偶函数就能秒判}}
\begin{center}\begin{tikzpicture}[font=\small,>=stealth,line join=round]
\draw[gray!40,rounded corners=3pt] (0,0) rectangle (14.2,8.3);\node at (7.1,7.700000000000001) {例4：余弦放在里面，整体自然左右同高};
\begin{scope}[shift={(4.1,2.4166666666666665)},x=0.6818181818181818cm,y=1.5333333333333332cm]\draw[->,gray] (-4.4,0)--(4.4,0) node[right]{$x$};\draw[->,gray] (0,-0.5)--(0,2.5) node[above]{$y$};\draw[blue,thick,solid,domain=-4:4,samples=180] plot(\x,{pow(2,cos(deg(\x)))});\end{scope}\node[anchor=west,align=left,text width=5.0cm] at (8.4,4.0) {$F=2^{\cos x}$\\[10pt]$\cos(-x)=\cos x$\\[10pt]同一个指数交给外层\\[10pt]$F(-x)=F(x)$};
\node[text width=12.6cm,align=center] at (7.1,.65) {曲线由原式计算；判断不必先展开整个指数函数。};
\end{tikzpicture}\end{center}


\textbf{题目：} 判断 $F(x)=2^{\cos x}$。\\
\textbf{第一步：} $g(x)=\cos x$，$f(u)=2^u$。因为 $\cos(-x)=\cos x$，内层是偶函数。\\
\textbf{第二步：} 指数函数 $2^u$ 不是偶函数，但这不影响结论：
\[
F(-x)=2^{\cos(-x)}=2^{\cos x}=F(x).
\]
\textbf{结果：} $F$ 为偶函数。以后看到“偶函数放在里面”，先不要急着计算整个大式子。

\medskip

\end{examplebox}
\clearpage
\begin{examplebox}
\textbf{\textcolor{CExample}{例 5：内奇外奇——正弦套立方}}
\begin{center}\begin{tikzpicture}[font=\small,>=stealth,line join=round]
\draw[gray!40,rounded corners=3pt] (0,0) rectangle (14.2,8.3);\node at (7.1,7.700000000000001) {例5：立方与正弦都把相反数保留下来};
\begin{scope}[shift={(4.1,3.9499999999999997)},x=1.5cm,y=1.6428571428571428cm]\draw[->,gray] (-2,0)--(2,0) node[right]{$x$};\draw[->,gray] (0,-1.4)--(0,1.4) node[above]{$y$};\draw[blue,thick,solid,domain=-1.8:1.8,samples=180] plot(\x,{sin(deg((\x)^3))});\end{scope}\node[anchor=west,align=left,text width=5.0cm] at (8.4,4.0) {$F=\sin(x^3)$\\[10pt]内层：$(-x)^3=-x^3$\\[10pt]外层：$\sin(-u)=-\sin u$\\[10pt]整体为奇函数};
\node[text width=12.6cm,align=center] at (7.1,.65) {起伏由真实公式决定；奇偶性关心原点对称，不要求单调。};
\end{tikzpicture}\end{center}


\textbf{题目：} 判断 $F(x)=\sin(x^3)$。\\
\textbf{第一步：} $g(x)=x^3$ 是奇函数，$f(u)=\sin u$ 也是奇函数。\\
\textbf{第二步：}
\[
F(-x)=\sin((-x)^3)=\sin(-x^3)=-\sin(x^3)=-F(x).
\]
\textbf{结果：} $F$ 为奇函数。这里的两次负号变换分别来自里面的立方与外面的正弦。

\medskip

\end{examplebox}
\clearpage
\begin{examplebox}
\textbf{\textcolor{CExample}{例 6：外层只有一半的输入范围，也能判偶}}
\begin{center}\begin{tikzpicture}[font=\small,>=stealth,line join=round]
\draw[gray!40,rounded corners=3pt] (0,0) rectangle (14.2,8.3);\node at (7.1,7.700000000000001) {例6：外层只需接受内层实际给出的数};
\begin{scope}[shift={(4.1,2.2716216216216214)},x=1.0714285714285714cm,y=1.243243243243243cm]\draw[->,gray] (-2.8,0)--(2.8,0) node[right]{$x$};\draw[->,gray] (0,-0.5)--(0,3.2) node[above]{$y$};\draw[blue,thick,solid,domain=-2.5:2.5,samples=180] plot(\x,{sqrt((\x)^2+1)});\fill[orange!85!black] (0,1) circle(2.2pt);\end{scope}\node[anchor=west,align=left,text width=5.0cm] at (8.4,4.0) {$F=\sqrt{x^2+1}$\\[10pt]内层实际输出 $[1,+\infty)$\\[10pt]全部能开平方\\[10pt]整体在 $\mathbb R$ 上为偶函数};
\node[text width=12.6cm,align=center] at (7.1,.65) {外层完整定义域不必左右对称；本题只使用非负的那部分输入。};
\end{tikzpicture}\end{center}


\textbf{题目：} 判断 $F(x)=\sqrt{x^2+1}$。\\
\textbf{第一步：} $g(x)=x^2+1$ 是偶函数，对一切实数 $x$，$g(x)\geqslant1$，所以能代入 $f(u)=\sqrt{u}$。\\
\textbf{第二步：}
\[
F(-x)=\sqrt{(-x)^2+1}=\sqrt{x^2+1}=F(x).
\]
\textbf{结果：} $F$ 为偶函数。外层只接受非负数，\textbf{根本不需要讨论外层是不是偶函数}。

\medskip

\end{examplebox}
\clearpage
\begin{examplebox}
\textbf{\textcolor{CExample}{例 7：反例——外层为偶，整体未必为偶}}
\begin{center}\begin{tikzpicture}[font=\small,>=stealth,line join=round]
\draw[gray!40,rounded corners=3pt] (0,0) rectangle (14.2,8.3);\node at (7.1,7.700000000000001) {外层平方为偶，整体却偏离了纵轴};
\begin{scope}[shift={(5.1,2.068181818181818)},x=1.3333333333333333cm,y=0.8363636363636363cm]\draw[->,gray] (-3,0)--(1.5,0) node[right]{$x$};\draw[->,gray] (0,-0.5)--(0,5) node[above]{$y$};\draw[blue,thick,solid,domain=-2.8:1.2,samples=180] plot(\x,{(\x+1)^2});\fill[orange!85!black] (-1,0) circle(2.2pt);\fill[orange!85!black] (1,4) circle(2.2pt);\end{scope}\node[anchor=west,align=left,text width=5.0cm] at (8.4,4.0) {$F=(x+1)^2$\\[10pt]内层 $g=x+1$ 非奇非偶\\[10pt]$F(1)=4,\ F(-1)=0$\\[10pt]对称轴 $x=-1$，不是 $y$ 轴};
\node[text width=12.6cm,align=center] at (7.1,.65) {外层只会把互为相反数的输入算成相同结果；这里的2和0并非相反数。};
\end{tikzpicture}\end{center}


\textbf{题目：} $f(u)=u^2$ 是偶函数，判断 $F(x)=f(x+1)$。\\
\textbf{拆成两层：} $g(x)=x+1$，不是奇函数，也不是偶函数。\\
\textbf{用 $1$ 和 $-1$ 检查：}
\[
F(1)=(1+1)^2=4,\qquad F(-1)=(-1+1)^2=0.
\]
$F(-1)$ 既不等于 $F(1)$，也不等于 $-F(1)$。\textbf{结果：} 既非奇，也非偶。图像其实关于 $x=-1$ 对称，并不关于 $y$ 轴对称。

\medskip

\end{examplebox}
\clearpage
\begin{examplebox}
\textbf{\textcolor{CExample}{例 8：不能倒推——整体为偶不能确定外层奇偶}}
\begin{center}\begin{tikzpicture}[font=\small,>=stealth,line join=round]
\draw[gray!40,rounded corners=3pt] (0,0) rectangle (14.2,8.3);\node at (7.1,7.700000000000001) {同样得到偶函数，外层的性质可以完全不同};
\begin{scope}[shift={(4.1,2.068181818181818)},x=2.0cm,y=0.41818181818181815cm]\draw[->,gray] (-1.5,0)--(1.5,0) node[right]{$x$};\draw[->,gray] (0,-1)--(0,10) node[above]{$y$};\draw[blue,thick,solid,domain=-1.2:1.2,samples=180] plot(\x,{(\x)^2+1});\draw[violet,thick,dashed,domain=-1.2:1.2,samples=180] plot(\x,{((\x)^2+1)^2+(\x)^2+1});\end{scope}\node[anchor=west,align=left,text width=5.0cm] at (8.4,4.0) {内层都为 $g=x^2+1$\\[10pt]蓝：外层 $f(u)=u$（奇）\\[10pt]紫：$f(u)=u^2+u$（非奇非偶）\\[10pt]两个整体都是偶函数};
\node[text width=12.6cm,align=center] at (7.1,.65) {两条实际曲线说明：仅从整体为偶，无法反推外层的奇偶性。};
\end{tikzpicture}\end{center}


\textbf{题目：} 已知 $F(x)=f(x^2+1)$ 是偶函数，能否判断 $f$ 的奇偶性？\\
\textbf{解释：} 因为 $x^2+1$ 已经把 $x$ 与 $-x$ 算成同一个数，外层只要能计算，最后就同高。\\
\textbf{举两个不同的外层：}
\[
\begin{array}{ll}
f(u)=u &\Longrightarrow F(x)=x^2+1,\\
f(u)=u^2+u &\Longrightarrow F(x)=(x^2+1)^2+(x^2+1).
\end{array}
\]
第一个 $f$ 是奇函数，第二个 $f$ 既非奇也非偶，但对应的 $F$ 都是偶函数。因此\textbf{不能只从整体为偶反推外层的奇偶性}。

\medskip

\end{examplebox}
\clearpage
\begin{examplebox}
\textbf{\textcolor{CExample}{例 9：让人眼前一亮——两层都不是奇偶函数，整体却是奇函数}}
\begin{center}\begin{tikzpicture}[font=\small,>=stealth,line join=round]
\draw[gray!40,rounded corners=3pt] (0,0) rectangle (14.2,8.3);\node at (7.1,7.700000000000001) {两次偏移抵消，整体也能出现奇偶性};
\begin{scope}[shift={(4.1,3.9499999999999997)},x=1.3043478260869565cm,y=1.0cm]\draw[->,gray] (-2.3,0)--(2.3,0) node[right]{$x$};\draw[->,gray] (0,-2.3)--(0,2.3) node[above]{$y$};\draw[blue,thick,solid,domain=-2:2,samples=180] plot(\x,{\x});\end{scope}\node[anchor=west,align=left,text width=5.0cm] at (8.4,4.0) {$g=x+1$，$f(u)=u-1$\\[10pt]两层本身都非奇非偶\\[10pt]$F=(x+1)-1=x$\\[10pt]整体却为奇函数};
\node[text width=12.6cm,align=center] at (7.1,.65) {口诀给出足够的条件，未穷尽所有复合函数的对称方式。};
\end{tikzpicture}\end{center}


\textbf{题目：} $g(x)=x+1$，$f(u)=u-1$，判断 $F(x)=f(g(x))$。\\
\textbf{先不要套口诀：} $g(-x)=-x+1$ 既非 $g(x)$，也非 $-g(x)$，所以口诀不能直接用。\\
\textbf{直接算：}
\[
F(x)=f(x+1)=(x+1)-1=x.
\]
因此 $F$ 是奇函数。\textbf{关键认识：} 口诀不是“有且只有”这些情况，它只是帮我们最快完成常见判断。

\medskip

\end{examplebox}
\clearpage
\begin{examplebox}
\textbf{\textcolor{CExample}{例 10：考场训练——条件不完整，先查定义域}}
\begin{center}\begin{tikzpicture}[font=\small,>=stealth,line join=round]
\draw[gray!40,rounded corners=3pt] (0,0) rectangle (14.2,8.3);\node at (7.1,7.700000000000001) {例10：定义域中间有缺口，左右仍成对};
\begin{scope}[shift={(4.1,2.307142857142857)},x=1.0714285714285714cm,y=1.3142857142857143cm]\draw[->,gray] (-2.8,0)--(2.8,0) node[right]{$x$};\draw[->,gray] (0,-0.5)--(0,3) node[above]{$y$};\draw[blue,thick,solid,domain=-2.5:-1,samples=180] plot(\x,{sqrt((\x)^2-1)});\draw[blue,thick,solid,domain=1:2.5,samples=180] plot(\x,{sqrt((\x)^2-1)});\fill[orange!85!black] (-1,0) circle(2.2pt);\fill[orange!85!black] (1,0) circle(2.2pt);\end{scope}\node[anchor=west,align=left,text width=5.0cm] at (8.4,4.0) {$F=\sqrt{x^2-1}$\\[10pt]$D=(-\infty,-1]\cup[1,+\infty)$\\[10pt]两支左右对应\\[10pt]整体仍是偶函数};
\node[text width=12.6cm,align=center] at (7.1,.65) {端点实心，因为根号允许被开方数为0；缺口本身不破坏定义域对称。};
\end{tikzpicture}\end{center}


\textbf{题目：} 判断 $F(x)=\sqrt{x^2-1}$ 的奇偶性。\\
\textbf{第一步：} 因为被开方数不能为负，$x^2-1\geqslant0$，所以允许 $x\leqslant-1$ 或 $x\geqslant1$，这个范围左右对称。\\
\textbf{第二步：} 里面 $g(x)=x^2-1$ 是偶函数，且 $g(x)\geqslant0$。\\
\textbf{第三步：}
\[
F(-x)=\sqrt{(-x)^2-1}=\sqrt{x^2-1}=F(x).
\]
\textbf{结果：} $F$ 是偶函数。\textbf{提醒：} 定义域不一定要是全体实数，只要它左右成对即可。
\end{examplebox}
\clearpage
\begin{examplebox}
\textbf{\textcolor{CExample}{例 11：负的平方作指数，仍是内偶则偶}}
\begin{center}
\begin{tikzpicture}[font=\small,>=stealth,line join=round]
\draw[gray!40,rounded corners=3pt] (0,0) rectangle (14.2,7.8);
\node at (7.1,7.2) {$a=2$：正负输入先得到同一个指数};
\begin{scope}[shift={(4.0,2.155172413793103)},x=1.2 cm,y=3.03448275862069 cm]
\draw[->,gray] (-2.5,0)--(2.5,0) node[right] {$x$};
\draw[->,gray] (0,-0.15)--(0,1.3) node[above] {$y$};
\draw[blue,thick,domain=-2.5:2.5,samples=201] plot (\x,{pow(2,-(\x)^2)});
\draw[orange!80!black,dashed] (-1,.5)--(1,.5);\fill[orange!80!black] (-1,.5) circle(2pt);\fill[orange!80!black] (1,.5) circle(2pt);
\end{scope}
\node[anchor=west,align=left,text width=5.1cm] at (8.1,4.0) {$F(x)=2^{-x^2}$\\[12pt]$-(-x)^2=-x^2$\\[12pt]$F(-1)=F(1)=\dfrac12$\\[12pt]内层相同，外层结果相同。};
\end{tikzpicture}
\end{center}

设 $a>0$ 且 $a\ne1$，判断 $F(x)=a^{-x^2}$ 的奇偶性。

定义域为 $\mathbb R$。内层 $g(x)=-x^2$ 满足
$g(-x)=-(-x)^2=-x^2=g(x)$，为偶函数；外层 $f(u)=a^u$ 在这些输入处都有定义。
因此
\[
F(-x)=a^{-(-x)^2}=a^{-x^2}=F(x),
\]
故 $F$ 为\textbf{偶函数}。回验：$F(-1)=F(1)=a^{-1}$。指数前的负号不会改变内层的偶性；图中只取 $a=2$ 示范，证明适用于两类底数。
\end{examplebox}
\clearpage
\begin{examplebox}
\textbf{\textcolor{CExample}{例 12：对数套绝对值，先排除零}}
\begin{center}
\begin{tikzpicture}[font=\small,>=stealth,line join=round]
\draw[gray!40,rounded corners=3pt] (0,0) rectangle (14.2,7.8);
\node at (7.1,7.2) {$a=2$：取绝对值后，两侧得到相同真数};
\begin{scope}[shift={(4.0,4.633333333333334)},x=1.0 cm,y=0.7333333333333334 cm]
\draw[->,gray] (-3,0)--(3,0) node[right] {$x$};
\draw[->,gray] (0,-4)--(0,2) node[above] {$y$};
\draw[blue,thick,domain=-3:-0.0625,samples=201] plot (\x,{ln(abs(\x))/ln(2)});
\draw[blue,thick,domain=0.0625:3,samples=201] plot (\x,{ln(abs(\x))/ln(2)});
\fill[orange!80!black] (-2,1) circle(2pt);\fill[orange!80!black] (2,1) circle(2pt);
\end{scope}
\node[anchor=west,align=left,text width=5.1cm] at (8.1,4.0) {$G(x)=\log_2|x|$\\[12pt]$G(-2)=G(2)=1$\\[12pt]$x\ne0$，定义域成对。\\[10pt]外层对数只接收正数，\\不必具有奇偶性。};
\end{tikzpicture}
\end{center}

设 $a>0$ 且 $a\ne1$，判断 $G(x)=\log_a|x|$ 的奇偶性。

先要求真数 $|x|>0$，得定义域 $\mathbb R\setminus\{0\}$，关于原点对称。
在这个范围内，内层 $g(x)=|x|$ 为偶函数，且始终为正，可以代入外层 $f(u)=\log_a u$。于是
\[
G(-x)=\log_a|-x|=\log_a|x|=G(x).
\]
故 $G$ 为\textbf{偶函数}。回验：$G(-1)=G(1)=0$，但 $G(0)$ 无定义。
这里用的是“内偶则偶”，不要求外层对数函数自身有奇偶性。

\textbf{关联阅读：} F083《常见指数、对数型函数的奇偶性速判》；$a^{|x|}$ 的翻折见 F047。
\end{examplebox}
